\section{Scattering Theory }
\subsection{Spherical Waves}
A spherical electromagnetic wave is of the following form:

\begin{equation}
	E(\mathbf{r}, t)=\frac{E_{0}}{\lvert \mathbf{r} \rvert }\exp(\mathrm{i}(k \lvert
	\mathbf{r}\rvert - \omega t))
\end{equation}

\subsection{Planar Approximation}

If the wave is observed in a small region $\mathcal{M}$ far away from its origin,
it can be approximated as a planar wave. Let the center of $\mathcal{M}$ be at
$\mathbf{R}$, and $\mathbf{r}\in \mathcal{M}$. By expressing $\mathbf{r}$
relative to $\mathbf{R}$ using the (small) distance vector
$\mathbf{r}-\mathbf{R}=\Delta{\mathbf{r}}$, an approximation for the magnitude can
be made:

\begin{align}
	\lvert \mathbf{r}\rvert & =\lvert \mathbf{R}+\Delta \mathbf{r}\rvert =\sqrt{ R^{2}+2\mathbf{R}\cdot\Delta \mathbf{r}+\Delta r^{2}}                         \\
	                        & \simeq R \sqrt{ 1+\frac{2\mathbf{R}\cdot\Delta \mathbf{r}}{R^{2}}}                                                               \\
	                        & \simeq R\cdot\left( 1+\frac{\mathbf{R}\cdot\Delta \mathbf{r}}{R^{2}}\right) =\hat{\mathbf{R}}\cdot(\mathbf{R}+\Delta \mathbf{r})
\end{align}
Defining $\mathbf{k}=k \hat{\mathbf{R}}$ and substituting the above relation into
the spherical electromagnetic wave equation directly results in a planar wave:

\begin{align}
	E(\mathbf{R}+\Delta\mathbf{r}, t) & =\frac{E_{0}}{\lvert \mathbf{R}+\Delta\mathbf{r} \rvert }\exp(\mathrm{i}(k \lvert \mathbf{R}+\Delta\mathbf{r}\rvert - \omega t)) \\
	                                  & \simeq \frac{E_{0}}{\lvert \mathbf{R} \rvert }\exp(\mathrm{i}(\mathbf{k}\cdot(\mathbf{R}+\Delta \mathbf{r})-\omega t))
\end{align}

\subsection{X-Ray Scattering Intensity}

A suitable framework to describe the interaction between light and solids is the kinematic
scattering theory. Assume an X-ray source emitting spherical waves far away from
the solid. The material center is located at $\mathbf{R}$. Every point inside
the material can be described by $\mathbf{r}=\mathbf{R}+\Delta \mathbf{r}$ with
$\Delta \mathbf{r}\ll \mathbf{R}$. As discussed in Spherical Wave, the wave
observed by the material is approximately planar. Defining $\mathbf{k}=k \hat{\mathbf{R}}$
and $k=2\pi /\lambda$ ($\lambda$ is the X-ray wavelength) leads to:

\begin{equation}
	E_{\mathrm{in}}(\mathbf{R}+\Delta\mathbf{r}, t) \simeq \frac{E_{0}}{\lvert
	\mathbf{R} \rvert }\exp\left(\mathrm{i}\left(\mathbf{k}_{\mathrm{in}}\cdot(\mathbf{R}
	+\Delta \mathbf{r})-\omega t\right)\right)
\end{equation}
The position of the material relative to the X-ray source dictates the direction of
the planar wave vector. Its magnitude is defined via $\lambda$. Incoming radiation
induces harmonic oscillations of shell electrons inside the solid. This leads to
spherical electromagnetic emission. The process is called Thomson scattering. It
is an elastic scattering mechanism that implies
$\lvert \mathbf{k}_{\mathrm{out}}\rvert=\lvert \mathbf{k}_{\mathrm{in}}\rvert$. The
outgoing spherical wave originating from a single point $\mathbf{P}=\mathbf{R}+\Delta
\mathbf{r}$ with scattering strength $f$ can be described as:

\begin{equation}
	E_{\mathrm{out}}(\mathbf{r},t)_{\mathbf{P}}=f E_{\mathrm{in}}(\mathbf{P}, t) \frac{\exp(\mathrm{i}k_{\mathrm{out}}\lvert
	\mathbf{r}-\mathbf{P}\rvert )}{\lvert \mathbf{r}-\mathbf{P} \rvert }
\end{equation}
Assume that an X-ray detector is located far away from the sample at position $\mathbf{D}
=\mathbf{R}+\mathbf{R}'$. The electric field at the detector's position is given
by:

\begin{equation}
	E_{\mathrm{out}}(\mathbf{D}, t)_{\mathbf{P}}=f E_{\mathrm{in}}(\mathbf{P},t) \frac{\exp(\mathrm{i}k_{\mathrm{out}}\lvert
	\mathbf{R}' -\Delta \mathbf{r}\rvert )}{\lvert \mathbf{R}'-\Delta \mathbf{r} \rvert
	}
\end{equation}
Since $\mathbf{R}'\gg \Delta \mathbf{r}$, the spherical wave can be approximated by
a planar wave using $\mathbf{k}_{\mathrm{out}}=k_{\mathrm{out}}\hat{\mathbf{R}}'$.
By setting the position of the detector relative to the sample, the direction of $\mathbf{k}
_{\mathrm{out}}$ is defined. Substituting $E_{\mathrm{in}}$ gives:

\begin{strip}
	\begin{align*}
		E_{\mathrm{out}}(\mathbf{D}, t)_{\mathbf{P}} & =\frac{f E_{0}}{\lvert \mathbf{R} \rvert \lvert \mathbf{R}' \rvert }\exp\left(\mathrm{i}\left(\mathbf{k}_{\mathrm{in}}\cdot(\mathbf{R}+\Delta \mathbf{r})-\omega t\right)\right) \exp\left(\mathrm{i}\mathbf{k}_{\mathrm{out}}\cdot(\mathbf{R}'-\Delta \mathbf{r})\right)                                                                                              \\
		                                             & =\frac{f}{\lvert \mathbf{R}' \rvert }\cdot \underbrace{ \frac{E_{0}}{\lvert \mathbf{R} \rvert } \exp\left(\mathrm{i}\left(\mathbf{k}_{\mathrm{in}}\cdot\mathbf{R}+\mathbf{k}_{\mathrm{out}}\cdot\mathbf{R}'\right)-\mathrm{i}\omega t\right) }_{ =E_{0}' }\cdot \exp\left(\mathrm{i}(\mathbf{k}_{\mathrm{in}}-\mathbf{k}_{\mathrm{out}})\cdot \Delta \mathbf{r}\right) \\
		                                             & =f \frac{E_{0}'}{\lvert \mathbf{R}' \rvert }\exp\left(\mathrm{i}(\mathbf{k}_{\mathrm{in}}-\mathbf{k}_{\mathrm{out}})\cdot\Delta \mathbf{r}\right)
	\end{align*}
\end{strip}
The expression above defines the electromagnetic response of a single atom. To find
the total material response, one needs to integrate this quantity with respect
to $\mathbf{P}$ over the crystal volume. Integration with respect to $\mathbf{P}$
is identical to integration with respect to $\Delta \mathbf{r}$. Since the
scattering is no longer originating from a single point but rather from a
density of electron scatterers, $f$ is replaced by the electron density $n(\Delta
\mathbf{r})$, resulting in:

\begin{strip}
	\begin{equation}
		E_{\mathrm{out}}(\mathbf{D}, t) \propto \frac{E_{0}'(\mathbf{R}, \mathbf{R}',
		t)}{\lvert \mathbf{R} '\rvert }\int_{V}n(\Delta \mathbf{r}) \exp\left(\mathrm{i}
		(\mathbf{k}_{\mathrm{in}}-\mathbf{k}_{\mathrm{out}})\cdot\Delta \mathbf{r}\right
		) \, \mathrm{d}V(\Delta \mathbf{r})
	\end{equation}
\end{strip}
By defining the scattering vector
$\mathbf{q}=\mathbf{k}_{\mathrm{out}}-\mathbf{k}_{\mathrm{in}}$, which depends only
on the experimental geometry, the intensity is expressed as:

\begin{equation}
	I(\mathbf{q}) \propto \frac{\lvert \langle E_{0}'(\mathbf{R}, \mathbf{R}') \rangle_{t}\rvert^{2}}{\lvert
	\mathbf{R}' \rvert^{2}}\left\lvert \int_{V}n(\mathbf{r}) \exp(-\mathrm{i}\mathbf{q}
	\cdot \mathbf{r}) \, \mathrm{d}V(\mathbf{r}) \right\rvert^{2}
\end{equation}
The intensity is proportional to the modulus of the Fourier transform of the
electron density in the crystal.

\subsection{Scattering Intensity for Periodic Structures}

The scattering intensity of an electron density far away from the electromagnetic
source and detector was derived in X-Ray Scattering Intensity and is given by:

\begin{equation}
	I(\mathbf{q}) \propto \left\lvert \int_{V}n(\mathbf{r}) \exp(-\mathrm{i}\mathbf{q}
	\cdot \mathbf{r}) \, \mathrm{d}V(\mathbf{r}) \right\rvert^{2}
\end{equation}
If the electron density is a periodic function, which is true for electrons in a
solid, it can be represented using the Reciprocal Lattice:

\begin{equation}
	I(\mathbf{q}) \propto \left\lvert \sum_{\mathbf{G}}n_{\mathbf{G}}\int_{V}\exp(\mathrm{i}
	(\mathbf{G}-\mathbf{q})\cdot\mathbf{r}) \, \mathrm{d}V(\mathbf{r}) \right\rvert
	^{2}
\end{equation}
The integral evaluates to:

\begin{align}
	\int_{V}\exp(\mathrm{i}(\mathbf{G}-\mathbf{q})\cdot\mathbf{r}) \, \mathrm{d}V(\mathbf{r}) & =\prod_{j=1}^{3}\int_{-l_{j} /2}^{l_{j} /2}\exp(\mathrm{i}(G_{j}-q_{j})x_{j})\, \mathrm{d}x_{j}                        \\
	                                                                                          & =\prod_{j=1}^{3}\left[ \frac{\exp(\mathrm{i}(G_{j}-q_{j})x_{j})}{\mathrm{i}(G_{j}-q_{j})}\right]^{l_{j} /2}_{-l_{j}/2} \\
	                                                                                          & =\prod_{j=1}^{3}\frac{\sin\left( l_{j}\frac{1}{2}(G_{j}-q_{j}) \right)}{\frac{1}{2}(G_{j}-q_{j})}                      \\
	                                                                                          & \simeq l_{x}l_{y}l_{z}\delta(\mathbf{q}, \mathbf{G})=V\delta(\mathbf{q}, \mathbf{G})
\end{align}
Substituting this back into the intensity expression yields the Laue condition:

\begin{align}
	I(\mathbf{q}) & \propto\left\lvert \sum_{\mathbf{G}}n_{\mathbf{G}}\int_{V}\exp(\mathrm{i}(\mathbf{G}-\mathbf{q})\cdot\mathbf{r}) \, \mathrm{d}V(\mathbf{r}) \right\rvert^{2}                                                             \\
	              & =\left\lvert \sum_{\mathbf{G}}n_{\mathbf{G}}V \delta(\mathbf{q}, \mathbf{G}) \right\rvert^{2}=\left\{ \begin{matrix}\lvert n_{\mathbf{q}}\rvert^{2}V^{2}&\mathbf{q}\in \mathcal{G}\\ 0&\mathrm{else}\end{matrix} \right.
\end{align}
Reflexes can only be observed if $\mathbf{q}=\mathbf{G}$. The scattering is
proportional to $V^{2}$ and therefore proportional to $N^{2}$, with $N$ being the
number of lattice points.

\subsection{Structure Factor}

\subsection{Derivation}

The Laue condition, see Scattering intensity for periodic structures, yields:

\begin{equation}
	I(\mathbf{q})\propto \left\{
	\begin{matrix}
		\lvert n_{\mathbf{q}}\rvert^{2}V^{2} & \mathbf{q}\in \mathcal{G} \\
		0                                    & \mathrm{else}
	\end{matrix}
	\right.
\end{equation}
As stated earlier, reflexes can only occur for $\mathbf{q}\in \mathcal{G}$. However,
if $n_{\mathbf{q}}=0$, no intensity is measured. These reflexes are called
forbidden. To determine when this happens, use the Fourier transform of the
electron density:

\begin{equation}
	n(\mathbf{q}) =\mathcal{F}^{-1}[n(\mathbf{r})] =\frac{1}{V_{\mathrm{uc}}}\int_{{\mathrm{uc}}}
	n(\mathbf{r}) \exp(-\mathrm{i}\mathbf{q}\cdot \mathbf{r}) \, \mathrm{d}V(\mathbf{r}
	)
\end{equation}
Due to the periodicity of the electron density, the region of integration reduces
to the unit cell with volume $V_{\mathrm{uc}}$. Since the main contribution to
the scattered X-rays originates from core electrons, where delocalization effects
can be neglected, the electron density can be approximated as a sum of atomic
densities $n_{\alpha}(\boldsymbol{\rho})$ within the unit cell:

\begin{align}
	n(\mathbf{r}) & = \sum_{\alpha}n_{\alpha}(\mathbf{r}-\mathbf{r}_{\alpha})                                                                                    \\
	              & = \sum_{\alpha}\int n_{\alpha}(\boldsymbol{\rho}) \delta(\mathbf{r}-\mathbf{r}_{\alpha}-\boldsymbol{\rho}) \, \mathrm{d}V(\boldsymbol{\rho}) \\
	              & =\sum_{\alpha}n_{\alpha}(\mathbf{r}) * \delta(\mathbf{r}-\mathbf{r}_{\alpha})
\end{align}
The last term is a sum over all basis atoms' electron densities convolved with the
delta distribution centered at $\mathbf{r}_{\alpha}$. Substituting this back
into the Fourier transform yields:

\begin{align}
	n(\mathbf{q}) & =\sum_{\alpha}\mathcal{F}^{-1}[\delta(\mathbf{r}-\mathbf{r}_{\alpha})] \cdot \mathcal{F}^{-1}[n_{\alpha}(\mathbf{r})]                                                                                                                                         \\
	              & =\frac{1}{V_{\mathrm{uc}}}\sum_{\alpha}\exp(-\mathrm{i}\mathbf{q}\cdot \mathbf{r}_{\alpha}) \underbrace{ \int n_{\alpha}(\boldsymbol{\rho}) \exp(-\mathrm{i} \mathbf{q}\cdot\boldsymbol{\rho}) \, \mathrm{d}V(\boldsymbol{\rho}) }_{ f_{\alpha}(\mathbf{q}) }
\end{align}
$f_{\alpha}(\mathbf{q})$ is the Fourier transform of the atomic electron density
and is also called the atomic scattering factor. The quantity $n(\mathbf{q})$ is also
known as the structure factor $S_{hkl}$ by setting $\mathbf{q}_{hkl}=h \mathbf{b}
_{1}+k\mathbf{b}_{2}+l\mathbf{b}_{3}$:

\begin{equation}
	S_{hkl}=\sum_{\alpha}f_{\alpha}(\mathbf{q}_{hkl}) \exp(-\mathrm{i}\mathbf{q}_{hkl}
	\cdot \mathbf{r}_{\alpha})
\end{equation}
Representing $\mathbf{r}_{\alpha}=u_{\alpha}\mathbf{a}_{1}+v_{\alpha}\mathbf{a}_{2}
+w_{\alpha}\mathbf{a}_{3}$ with $u_{\alpha}, v_{\alpha}, w_{\alpha}\in[0, 1)$,
the structure factor can be expressed as

\begin{equation}
	S_{hkl}=\sum_{\alpha}f_{\alpha}(\mathbf{q}_{hkl}) \exp(-2\pi \mathrm{i}(hu_{\alpha}
	+kv_{\alpha}+l w_{\alpha}))
\end{equation}
\subsection{Example}

The ideal wurtzite lattice contains four basis atoms. Oxygen atoms are located at
$(000)$ and $\left( \frac{1}{3}\frac{2}{3}\frac{1}{2}\right)$, zinc atoms at
$\left ( 0 0 \frac{3}{8}\right)$ and
$\left( \frac{1}{3}\frac{2}{3}\frac{7}{8}\right)$. By assuming constant scattering
factors for zinc and oxygen, respectively, the structure factor simplifies to:

\begin{align}
	S_{hkl} & =f_{\ce{O}}\left[ 1+\exp\left( -2\pi \mathrm{i}\left( \frac{h}{3}+\frac{2k}{3}+\frac{l}{2}\right) \right) \right]                                                                  \\
	        & +f_{\ce{Zn}}\left[ \exp\left( -2\pi \mathrm{i}\left( \frac{3l}{8}\right) \right) + \exp\left( -2\pi \mathrm{i}\left( \frac{h}{3}+\frac{2k}{3}+\frac{7}{8}l \right) \right) \right]
\end{align}
Analyzing the $(00L)$ directions reveals forbidden reflexes:

\begin{equation}
	S_{00L}=2 \left( f_{\ce{O}}+f_{\ce{Zn}}\exp\left( -\frac{3\pi}{4}\mathrm{i}l \right
	) \right) \mathbb{1}_{l \text{ even}}
\end{equation}

\subsection{Bragg Equation}

The Laue condition from Scattering intensity for periodic structures states that reflexes
can only be observed if $\mathbf{q}$ matches a reciprocal space vector,
therefore $\mathbf{q}=z \mathbf{G}_{hkl}$ with $z \in\mathbb{Z}$. A visual explanation
of this result is obtained using the lattice spacing equation from Reciprocal Lattice:
$d_{hkl}={2\pi}/{\lvert \mathbf{G}_{hkl} \rvert }$. Representing $\lvert \mathbf{q}
\rvert$ using the diffractometer angles from XRD Coordinate Systems yields the \textbf{Bragg
equation}:

\begin{align}
	d_{hkl}  & =\frac{2\pi n}{2\sin\theta\cdot2\pi /\lambda} \\
	n\lambda & =2d_{hkl}\sin( \theta)
\end{align}
The lattice planes partially reflect the incident wave. If the path difference for
waves reflected by adjacent lattice planes is an integer multiple of the wavelength,
constructive interference is expected.

\subsection{Maximum Error Estimation}

Performing maximum error estimation of the Bragg equation results in:

\begin{equation}
	\frac{\Delta\lambda}{\lambda}=\cot(\theta) \Delta\theta+\frac{\Delta d_{hkl}}{d_{hkl}}
\end{equation}
To determine the distance $d_{hkl}$ accurately, it is essential to use a highly
monochromatic incident beam and a detector with excellent angular resolution.

\subsection{Ewald Sphere}

The Laue condition, derived in Scattering intensity for periodic structures, states:

\begin{align}
	I(\mathbf{q}) & \propto \left\{ \begin{matrix}\lvert n_{\mathbf{q}}\rvert^{2}V^{2}&\mathbf{q}\in \mathcal{G}\\ 0&\mathrm{else}\end{matrix} \right.
\end{align}
The intensity is only nonzero if $\mathbf{q}$ is a reciprocal space vector. Each
intensity maximum can therefore be labeled with its corresponding reciprocal
lattice index $( hkl)$. Since the magnitude $\lvert \mathbf{k}_{\mathrm{in}}\rvert
=\lvert \mathbf{k}_{\mathrm{out}}\rvert$ is defined as $k=2 \pi /\lambda$, the
scattering vector satisfies

\begin{equation}
	\lvert \mathbf{q}\rvert =\lvert \mathbf{k}_{\mathrm{out}}-\mathbf{k}_{\mathrm{in}}
	\rvert \leq 2k
\end{equation}
In reciprocal space, this defines the \textbf{Ewald Sphere} for $\mathbf{q}$ with
radius $2k$. Observable reflexes correspond to reciprocal lattice points within
that sphere. Assuming that the laboratory system is aligned with the crystallographic
coordinate system, see XRD Coordinate Systems, leads to the diagram below: